\chapter{Conclusion and Future Work}

\section{Summary of Contributions}

\begin{table}[H]
\centering
\resizebox{0.7\textwidth}{!}{
\begin{tabular}{p{6cm}p{7cm}}
\toprule
\textbf{Contribution} & \textbf{Outcome} \\
\midrule
LLM extraction pipeline
  & Structured disruption events from 9 live sources;
    13 event types; mandatory location enforcement \\[4pt]
Temporal HGNN
  & City-wide confidence refinement; 3-output-head model;
    cascade road detection; attention-based explainability \\[4pt]
Bayesian probabilistic fusion
  & Per-road posterior disruption probabilities;
    time-of-day calibrated priors; HGNN blend \\[4pt]
Multimodal routing
  & 8 transport modes; full 5-line Kolkata Metro;
    BFS bus routing; OSMnx road/walk/bike graphs \\[4pt]
Weather risk assessment
  & Route-point sampling; Weather Severity Index (WSI)
    per candidate route \\[4pt]
FastAPI backend
  & 11 production endpoints; two-step design;
    2~s initial render regardless of LLM latency \\[4pt]
Frontend
  & Leaflet.js map; risk-coloured route lines;
    event markers; metro + bus overlays \\[4pt]
Database
  & Auto-migrating SQLite schema; HGNN feedback loop;\\
\bottomrule
\end{tabular}
}
\caption{Contributions made during the internship}
\end{table}

\section{Research Hypotheses -- Status}

\begin{table}[H]
\centering
\resizebox{0.7\textwidth}{!}{
\begin{tabular}{cp{7cm}p{4.5cm}}
\toprule
\textbf{ID} & \textbf{Hypothesis} & \textbf{Status} \\
\midrule
H1 & LLM signals improve disruption detection accuracy
   & Demonstrated qualitatively; ground-truth eval pending \\[4pt]
H2 & Risk-aware routing improves travel-time reliability
   & Layer~2 complete; Layer~3 (CVaR) in progress \\[4pt]
H3 & Explainable outputs enhance user trust
   & HGNN attention exposed via API; user study pending \\
\bottomrule
\end{tabular}
}
\caption{Hypothesis status at end of internship}
\end{table}

\section{Skills Acquired}

\begin{table}[H]
\centering
\begin{tabular}{ll}
\toprule
\textbf{Area} & \textbf{Specifics} \\
\midrule
LLM application dev    & LangChain LCEL, prompt engineering, JSON repair,
                         structured output \\
Graph neural networks  & Heterogeneous graph construction, HAN layers,
                         multi-task training (MSE + CrossEntropy) \\
Probabilistic ML       & Bayesian inference, likelihood ratios,
                         time-varying priors \\
Geospatial engineering & OSMnx, OpenStreetMap, GeoJSON, Leaflet.js,
                         Nominatim geocoding \\
Backend development    & FastAPI, Pydantic, SQLAlchemy, REST API design \\
Data engineering       & Multi-source ingestion, RSS/web scraping,
                         API rate-limit handling \\
\bottomrule
\end{tabular}
\caption{Technical skills developed}
\end{table}

\section{Future Work}

\begin{table}[H]
\centering
\resizebox{0.7\textwidth}{!}{
\begin{tabular}{cp{5.5cm}p{6cm}}
\toprule
\textbf{\#} & \textbf{Task} & \textbf{Details} \\
\midrule
1 & Layer 3 --- CVaR routing
  & Generalised edge cost combining travel time, risk, CO$_2$,
    and transfer penalty; CVaR optimisation for route selection \\[4pt]
2 & HGNN ground-truth training
  & Collect verified disruption labels from Kolkata Traffic Police;
    retrain severity and confidence heads on real data \\[4pt]
3 & GTFS-RT integration
  & Replace static metro/bus timetables with live GTFS-Realtime feeds \\[4pt]
4 & Natural language explanations
  & Layer~3 generates a plain-English justification per route
    recommendation (H3 full validation) \\[4pt]
5 & User trust study
  & Formal user experiment to validate H3 with Kolkata commuters \\[4pt]
6 & Multi-city deployment
  & Replace Kolkata OSM graph and transport datasets with those
    of another Indian metro for portability test \\
\bottomrule
\end{tabular}
}
\caption{Planned extensions}
\end{table}
